\documentclass[10pt]{article}
\usepackage[T1]{fontenc}
\usepackage[utf8]{inputenc}
\usepackage[french]{babel}
\usepackage{amsmath,amssymb,mathrsfs,mathpazo}
\usepackage{geometry}
\geometry{a4paper,margin=2cm}
\usepackage{graphicx,array,tabularx,tabularray,colortbl,multirow,multicol,booktabs,enumitem,pifont,ulem,fancybox,diagbox,cancel}
\usepackage{pstricks,pstricks-add,pst-plot,pst-tree,pst-node,pst-text,pst-func,pst-math,pst-bspline,pst-3dplot,pst-eucl,pst-solides3d}
\let\pscalculate\undefined
\usepackage{pst-bezier}
\usepackage{tikz,pgfplots,tkz-tab}
\usetikzlibrary{arrows,arrows.meta,patterns,calc,positioning,trees,decorations.markings,shapes.geometric,shapes.misc,shapes.arrows,intersections,decorations.pathmorphing,decorations.pathreplacing,backgrounds,shadows,babel}
\pgfplotsset{compat=1.18}
\usepgfplotslibrary{fillbetween}
\usepackage[np]{numprint}
\usepackage{listings,makecell,pas-tableur,twoopt}
\let\barre\undefined
\usepackage[pstricks]{tablvar}
\definecolor{darkgreen}{RGB}{0,100,0}
\providecommand{\red}{\color{red}}
\providecommand{\blue}{\color{blue}}
\providecommand{\green}{\color{green}}
\providecommand{\black}{\color{black}}
\providecommand{\white}{\color{white}}
\providecommand{\gray}{\color{gray}}
\providecommand{\cyan}{\color{cyan}}
\providecommand{\magenta}{\color{magenta}}
\providecommand{\yellow}{\color{yellow}}
\providecommand{\violet}{\color{violet}}
\providecommand{\orange}{\color{orange}}
\providecommand{\euro}{\texteuro}
\providecommand{\EUR}{\texteuro}
\providecommand{\ieme}{\textsuperscript{e}}
\providecommand{\degres}{\ensuremath{^\circ}}
\providecommand{\degree}{\ensuremath{^\circ}}
\providecommand{\vv}[1]{\overrightarrow{#1}}
\providecommand{\Coord}[1]{#1}
\providecommand{\decosix}{}
\providecommand{\starredbullet}{\ensuremath{\bullet}}
\pagestyle{empty}
\setlength{\parindent}{0pt}
\usepackage{adjustbox}\DeclareUnicodeCharacter{2236}{:}\DeclareUnicodeCharacter{2190}{\ensuremath{\leftarrow}}
\providecommand{\R}{\mathbb{R}}
\providecommand{\N}{\mathbb{N}}
\providecommand{\Z}{\mathbb{Z}}
\providecommand{\Q}{\mathbb{Q}}
\providecommand{\D}{\mathbb{D}}
\providecommand{\C}{\mathbb{C}}
\providecommand{\e}{{\mathrm{e}}}
\providecommand{\vect}[1]{\overrightarrow{#1}}
\providecommand{\vectt}[1]{\overrightarrow{\mathrm{#1}}}
\providecommand{\vvt}[1]{\overrightarrow{\mathrm{#1}}}
\providecommand{\np}[2][]{{#2~#1}}
\providecommand{\vv}[1]{{\overrightarrow{#1}}}
\renewcommand{\d}{{\mathrm{d}}}
\providecommand{\degres}{{^\circ}}
\providecommand{\degre}{{^\circ}}
\providecommand{\blue}{\color{blue}}
\providecommand{\red}{\color{red}}
\providecommand{\enskip}{\quad}
\providecommand{\llbracket}{\lbrack\!\lbrack}
\providecommand{\rrbracket}{\rbrack\!\rbrack}
\providecommand{\euro}{{\text{€}}}
\providecommand{\no}{{\text{n°}}}
\providecommand{\interval}[2]{{[#1;#2]}}
\providecommand{\intervFO}[2]{{[#1;#2[}}
\providecommand{\dotfill}{\ldots}

\providecommand{\ds}{\displaystyle}
\providecommand{\textding}[1]{#1}
\providecommand{\Oij}{\(\mathrm{O};\vec{\imath},\vec{\jmath}\)}
\providecommand{\Oijk}{\(\mathrm{O};\vec{\imath},\vec{\jmath},\vec{k}\)}
\providecommand{\Ouv}{\(\mathrm{O};\vec{u},\vec{v}\)}
\providecommand{\intervFF}[2]{\ensuremath{[#1;#2]}}
\renewcommand{\euro}{\ifmmode\text{\texteuro}\else\texteuro\fi}
\DeclareUnicodeCharacter{00B0}{\ifmmode^{\circ}\else\textdegree\fi}
\DeclareUnicodeCharacter{20AC}{\ifmmode\text{\texteuro}\else\texteuro\fi}
% Énoncés : Éducation nationale. Transcription et corrigés : APMEP, Denis Vergès et contributeurs. Découpage et classement : CoopMaths, Sébastien Lozano, Éric Elter et contributeurs. Adaptation de la présentation : Maths Explorer.
% CC BY-SA : https://coopmaths.fr/www/about/
% https://coopmaths.fr/alea/static/bac/2022/tex/bac_2022_05_sujet2_polynesie_1.tex
\begin{document}
\noindent\textit{APMEP / CoopMaths. CC BY-SA. Présentation adaptée par Maths Explorer.}\par\medskip
\medskip

\emph{Cet exercice est un questionnaire à choix multiples. \\
Pour chacune des six questions suivantes, une seule des quatre réponses proposées est exacte.\\
Une réponse fausse, une réponse multiple ou l'absence de réponse à une question ne rapporte ni n'enlève de point.\\
Pour répondre, indiquer sur la copie le numéro de la question et la lettre de la réponse choisie.\\
Aucune justification n'est demandée.}

\medskip

\begin{enumerate}
\item On considère la fonction $f$ définie et dérivable sur $]0~;~+\infty[$ par:
\[f(x) = x \ln(x) - x + 1.\]

Parmi les quatre expressions suivantes, laquelle est celle de la fonction dérivée de $f$?

\begin{center}
\begin{tabularx}{\linewidth}{|*{4}{X|}}\hline
\textbf{a.~~} $\ln (x)$&\textbf{b.~~}$\dfrac{1}{x} - 1$&\textbf{c.~~} $\ln (x) - 2$&\textbf{d.~~}$\ln (x) - 1$\rule[-3mm]{0mm}{9mm}\\ \hline
\end{tabularx}
\end{center}

\item On considère la fonction $g$ définie sur $]0~;~+\infty[$ par $g(x) = x^2[1 - \ln (x)]$.

 Parmi les quatre affirmations suivantes, laquelle est correcte ?
 
\begin{center}
\begin{tabularx}{\linewidth}{|*{4}{X|}}\hline
\textbf{a.~~}$\displaystyle\lim_{x \to 0} g(x) = +\infty$&\textbf{b.~~} 
$\displaystyle\lim_{x \to 0} g(x) = - \infty$&\textbf{c.~~} $\displaystyle\lim_{x \to 0} g(x) = 0$&\textbf{d.~~} La fonction $g$ n'admet pas de limite en 0.\\ \hline
\end{tabularx}
\end{center}

\item On considère la fonction $f$ définie sur $\R$ par $f(x) = x^3 - 0,9x^2 -0,1x$. Le nombre de solutions de l'équation $f(x) = 0$ sur $\R$ est :

\begin{center}
\begin{tabularx}{\linewidth}{|*{4}{X|}}\hline
\textbf{a.~~} $0$&\textbf{b.~~}$1$&\textbf{c.~~} $2$&\textbf{d.~~}$3$\\ \hline
\end{tabularx}
\end{center}

\item Si $H$ est une primitive d'une fonction $h$ définie et continue sur $\R$, 
et si $k$ est la fonction définie sur $\R$ par $k(x) = h(2x)$, 
alors, une primitive $K$ de $k$ est définie sur $\R$ par :

\begin{center}
\begin{tabularx}{\linewidth}{|*{4}{X|}}\hline
\textbf{a.~~} $K(x) =H(2x)$&\textbf{b.~~}$K(x) =2H(2x)$&\textbf{c.~~} $ K(x) =\dfrac{1}{2}H(2x)$&\textbf{d.~~}$K(x) =2H(x)$\rule[-3mm]{0mm}{9mm}\\ \hline
\end{tabularx}
\end{center}

\item L'équation réduite de la tangente au point d'abscisse 1 de la courbe de la fonction $f$ définie sur $\R$ par $f(x) = x\text{e}^x$ est:

\begin{center}
\begin{tabularx}{\linewidth}{|*{4}{X|}}\hline
\textbf{a.~~} $y = \text{e}x + \text{e}$&\textbf{b.~~}$y =2\text{e}x - \text{e}$&\textbf{c.~~} $y = 2\text{e}x +  \text{e}$&\textbf{d.~~}$y = \text{e}x$\\ \hline
\end{tabularx}
\end{center}

\item Les nombres entiers $n$ solutions de l'inéquation $(0,2)^n < 0,001$ sont tous les
nombres entiers $n$ tels que :

\begin{center}
\begin{tabularx}{\linewidth}{|*{4}{X|}}\hline
\textbf{a.~~} $n \leqslant 4$&\textbf{b.~~}$n\leqslant 5$&\textbf{c.~~} $n \geqslant 4$ &\textbf{d.~~}$n \geqslant 5$\\ \hline
\end{tabularx}
\end{center}

\end{enumerate}

\bigskip

\end{document}
